\documentclass[12pt]{article}

\usepackage[T1]{fontenc}
\usepackage[utf8]{inputenc}
\usepackage[romanian]{babel}
\usepackage{amsmath, amssymb}
\usepackage{geometry}
\geometry{a4paper, margin=2.5cm}

\title{Algoritmi de Sortare — Prezentare Explicativă}
\author{}
\date{}

\begin{document}

\maketitle

\section{Introducere}

Sortarea este una dintre cele mai fundamentale operații din informatică. 
Apare în algoritmi de căutare, baze de date, optimizare, grafuri, procesarea datelor
și în aproape orice aplicație practică. 

Există numeroase metode de sortare, fiecare cu avantaje, dezavantaje și domenii
în care excelează. În acest capitol prezentăm principalele sortări clasice,
explicând intuitiv modul lor de funcționare, proprietățile și complexitățile.

\section{Sortarea prin interschimbare (Bubble Sort)}

Bubble Sort este una dintre cele mai simple sortări. Ideea este intuitivă:
parcurgem vectorul de mai multe ori și comparăm elemente adiacente. Dacă sunt în
ordine greșită, le interschimbăm. Elementele mari ``urcă'' treptat spre final,
ca niște bule într-un lichid.

Algoritmul este stabil și funcționează in-place, dar este foarte lent. 
În cel mai bun caz (vector deja sortat) poate fi liniar dacă detectăm că nu s-au
făcut schimbări. În medie și în cel mai rău caz are complexitate $O(n^2)$.

Este folosit aproape exclusiv în scop educațional.

\section{Sortarea prin selecție (Selection Sort)}

Selection Sort funcționează prin alegerea repetată a celui mai mic element din
partea nesortată și mutarea lui la început. La fiecare pas, secvența sortată
crește cu un element.

Algoritmul este simplu, folosește foarte puține interschimbări (cel mult $n-1$),
dar necesită $O(n^2)$ comparații indiferent de situație. Nu este stabil, deoarece
minimul poate proveni dintr-o poziție care schimbă ordinea relativă a elementelor
egale.

Este util când costul comparațiilor este mic, dar costul schimburilor este mare.

\section{Sortarea prin inserție (Insertion Sort)}

Insertion Sort imită modul în care oamenii sortează cărțile în mână: luăm
elementele pe rând și le inserăm în poziția corectă în partea deja sortată.

Este stabil, funcționează in-place și este foarte eficient pentru vectori aproape
sortați, având complexitate $O(n)$ în cel mai bun caz. În medie și în cel mai rău
caz are complexitate $O(n^2)$.

Este folosit frecvent ca subrutină în algoritmi mai mari, deoarece este rapid pe
intrări mici.

\section{Counting Sort}

Counting Sort este o sortare fără comparații, construită pe o idee esențială:
\textbf{poziția finală a unui element într-un vector sortat este determinată de câte elemente sunt strict mai mici decât el}.

Pentru a afla această informație eficient, algoritmul parcurge trei etape:

\begin{enumerate}
    \item \textbf{Numără aparițiile fiecărei valori} din intervalul permis. 
    Acesta este doar un pas intermediar, nu scopul final.

    \item \textbf{Transformă numărătorile în prefix-sum}, astfel încât pentru fiecare valoare $v$ să știm:
    

\[
    C[v] = \text{numărul de elemente } \le v.
    \]


    Din această informație rezultă imediat:
    

\[
    \text{numărul de elemente strict mai mici decât } v = C[v-1].
    \]



    \item \textbf{Folosind aceste valori}, fiecare element își poate determina direct poziția finală în vectorul sortat.
\end{enumerate}

Counting Sort este \textbf{stabil}, foarte rapid și are complexitate:


\[
O(n + k),
\]


unde $k$ este valoarea maximă posibilă. Devine însă ineficient dacă $k$ este foarte mare, deoarece necesită memorie suplimentară proporțională cu $k$.

Este folosit frecvent ca etapă în Radix Sort, unde domeniul valorilor este controlat.

\section{Merge Sort (sortarea prin interclasare)}

Merge Sort este un algoritm clasic de tip \textit{divide et impera}. Ideea lui este elegantă și puternică:

\begin{enumerate}
    \item împarte vectorul în două jumătăți aproximativ egale;
    \item sortează recursiv fiecare jumătate;
    \item interclasează cele două jumătăți sortate într-un singur vector ordonat.
\end{enumerate}

Interclasarea este operația-cheie: două secvențe deja sortate pot fi combinate într-una singură în timp liniar.

Merge Sort este un algoritm:
\begin{itemize}
    \item \textbf{stabil};
    \item \textbf{foarte robust}, având complexitate garantată:
    

\[
    O(n \log n)
    \]


    în orice caz;
    \item \textbf{nu este in-place}, deoarece necesită memorie suplimentară $O(n)$.
\end{itemize}

Este preferat în aplicații unde stabilitatea este importantă sau când datele sunt prea mari pentru a încăpea în memorie și sortarea se face pe disc (sortare externă).

\section{Quick Sort}

Quick Sort este considerat în practică unul dintre cei mai rapizi algoritmi de
sortare. Alege un pivot, împarte vectorul în elemente mai mici și mai mari decât
pivotul, apoi sortează recursiv cele două părți.

În medie are complexitate $O(n \log n)$, dar în cel mai rău caz poate ajunge la
$O(n^2)$ dacă pivotul este ales prost. Există însă strategii de alegere a
pivotului care reduc semnificativ riscul cazului nefavorabil.

Este in-place, dar nu este stabil. În practică este adesea cel mai rapid algoritm
de sortare generală.

\section{Heap Sort}

Heap Sort folosește o structură de date numită heap (un arbore binar complet cu
proprietatea de heap). Algoritmul construiește un heap din vector, apoi extrage
repetat maximul (sau minimul) și îl mută la final.

Are complexitate $O(n \log n)$ în orice caz, este in-place și nu necesită memorie
suplimentară semnificativă. Nu este stabil.

Este folosit când avem nevoie de performanță garantată și memorie minimă.

\section{Sortarea folosind arbori de căutare (BST Sort)}

Sortarea prin arbori de căutare constă în inserarea tuturor elementelor într-un
arbore binar de căutare, apoi efectuarea unei parcurgeri in-order, care produce
elementele în ordine sortată.

Complexitatea depinde de forma arborelui:
\begin{itemize}
    \item $O(n \log n)$ în medie,
    \item $O(n^2)$ în cel mai rău caz (arbore degenerat).
\end{itemize}

Dacă folosim arbori echilibrați (AVL, Red-Black), complexitatea devine garantat
$O(n \log n)$.

Algoritmul nu este in-place și nu este stabil.

\section{Concluzii}

Există numeroase algoritmi de sortare, fiecare potrivit pentru situații diferite.
Algoritmii simpli (Bubble Sort, Selection Sort, Insertion Sort) sunt utili pentru
înțelegerea conceptelor de bază. Algoritmii avansați (Merge Sort, Quick Sort,
Heap Sort) sunt folosiți în practică datorită eficienței lor.

Sortările fără comparații (Counting Sort) pot fi extrem de rapide atunci când
domeniul valorilor este restrâns.

\end{document}
